% \iffalse meta-comment
%
% File: muling.dtx
% ----------------------------------------------------------
% Class:       muling
% Author:      निरंजन
% Description: A class file for the Department of
%              Linguistics University of Mumbai
% Repository:  https://git.gnu.org.ua/marathi.git
% Bugs:        https://puszcza.gnu.org.ua/bugs/?group=muling
% Licenses:    GPLv3+, GFDLv1.3+
% ----------------------------------------------------------
% LaTeX Package muling v0.5
% Copyright © 2020, 2021, 2023 निरंजन
%
% This program is free software: you can redistribute it
% and/or modify it under the terms of the GNU General Public
% License as published by the Free Software Foundation,
% either version 3 of the License, or (at your option) any
% later version.
%
% This program is distributed in the hope that it will be
% useful, but WITHOUT ANY WARRANTY; without even the implied
% warranty of MERCHANTABILITY or FITNESS FOR A PARTICULAR
% PURPOSE. See the GNU General Public License for more
% details.
%
% You should have received a copy of the GNU General Public
% License along with this program. If not, see
% <https://www.gnu.org/licenses/>.
% ---------------------------------------------------------
% arara: lualatex
% arara: lualatex
% arara: clean: {
% arara: -->      extensions: [
% arara: -->        log,aux,glo,idx,ist,
% arara: -->        out,toc,hd
% arara: -->      ]
% arara: -->    }
% ---------------------------------------------------------
% \fi
% \iffalse
%<*internal>
\iffalse
%</internal>
%<*readme>
-----------------------------------------------------------
Class:        muling
Author:       निरंजन
Version:      0.5 (2 April, 2023)
Description:  A class file for the Department of
              Linguistics, University of Mumbai
Repository:   https://git.gnu.org.ua/muling.git
Bug tracker:  https://puszcza.gnu.org.ua/bugs/?group=muling
License:      GPLv3+, GFDLv1.3+
-----------------------------------------------------------
%</readme>
%<*internal>
\fi
%</internal>
%<*driver>
\documentclass{l3doc}
\usepackage{xcolor}
\usepackage{leipzig}
\usepackage{fontspec}
\newfontfamily{\sho}[%
  Script          = {Devanagari},%
  Scale           = {0.8}%
]{Shobhika}
\usepackage{fontawesome5}
\usepackage{gfdl}
\colorlet{myblue}{blue!60!black}
\colorlet{myred}{red!60!black}
\colorlet{mygreen}{green!60!black}
\hypersetup{%
  unicode,%
  colorlinks,%
  linkcolor       = {myred},%
  citecolor       = {mygreen},%
  urlcolor        = {myblue},%
  pdftitle        = {The muling package},%
  pdfauthor       = {निरंजन},%
  pdfcreator      = {निरंजन},%
  pdfsubject      = {%
    A convenient LaTeX class for the masters students of
    linguistics at the department of linguistics, university
    of Mumbai%
  },%
  pdfkeywords     = {%
    LaTeX, Mumbai University, Thesis%
  }%
}
\gfdlcopyrightholders{%
  \texorpdfstring{{\sho निरंजन}}%
                 {निरंजन}%
}
\gfdlcopyrightdescription{%
  \texorpdfstring{The \LaTeX\ package \pkg{muling}}%
                 {The LaTeX package muling}%
}
\gfdlcopyrightableyears{2020, 2021, 2023}

\begin{document}
\DocInput{muling.dtx}
\end{document}
%</driver>
% \fi
% \title{The \cls{muling} class}
% \author{{\sho निरंजन}}
% \date{^^A
%   Version 0.5---2 April, 2023\\[1ex]^^A
%   {^^A
%     \small\faIcon{link}\quad
%     \url{https://ctan.org/pkg/muling}^^A
%   }\\^^A
%   {^^A
%     \small\faIcon{bug}\quad
%     \url{https://puszcza.gnu.org.ua/bugs/?group=muling}^^A
%   }^^A
% }
% 
% \maketitle
% 
% \begin{abstract}
%   This is a class-file written for the students of the
%   Department of Linguistics, University of Mumbai. It
%   uses the standard layout prescribed by the university,
%   loads some necessary packages for typesetting a thesis
%   in linguistics \textit{\&} adds the declaration needed
%   by the university from the submitter.
% \end{abstract}
% 
% \tableofcontents
% 
% \newpage
% 
% \begin{documentation}
% \section{Provisions}
% 
% With this class, a few essential packages are already
% loaded. If you have a full installation of \TeX\ Live
% 2023 (or later), you should be able to compile everything
% without any problem. This package has the following
% provisions which will make the thesis-writing easier.
% 
% \begin{enumerate}
% \item With this class, we load the |NewCM10-Regular.otf|
%   font which is a fork of Latin Modern with extended
%   support for Unicode IPA. Using this font enables us to
%   print IPA directly from the input text. An additional
%   advantage of using this font is that it matches with the
%   TIPA typeface exactly. So whether you use Unicode or
%   TIPA visually won't make any difference.
% \item The class pre-loads the following packages commonly
%   used in linguistics.
%   \begin{description}
%   \item[\pkg{leipzig}:] Used for inter-linear glossing.
%     Using |\Nom| will produce \Nom. While glossing, it is
%     advised that you use these macros so that all of them
%     get printed correctly with the |\printglossary|
%     command pre-loaded by the class. Of course you will be
%     required to compile the code with the |makeglossaries|
%     program.
%   \item[\pkg{biblatex}:] Used for citations. The |apa|
%     style, required by the department, is used by default.
%   \item[\pkg{forest}:] Used for syntactic trees.
%   \item[\pkg{expex}:] Used for finest typesetting of
%     linguistic examples.
%   \item[\pkg{csquotes}:] For quoting text. One can use
%   \cs{enquote}\marg{some text} to produce
%   \emph{\enquote{some text}}.
%   \end{description}
% \end{enumerate}
% 
% It is advised that you use the following order on a bash
% terminal for compiling your code:
% 
% \medskip
% \begin{minipage}{\linewidth}
% \begin{verbatim}
% lualatex document         # for Unicode text
% biber document            # for citations
% makeglossaries document   # for glossaries
% lualatex document
% lualatex document
% \end{verbatim}
% \end{minipage}
%
% \medskip
% If you use \pkg{arara}, then you can use the following
% code somewhere in your document and just compile the code
% with \pkg{arara}.
% 
% \medskip
% \begin{minipage}{\linewidth}
% \begin{verbatim}
% % arara: lualatex
% % arara: biber
% % arara: makeglossaries
% % arara: lualatex
% % arara: lualatex
% \end{verbatim}
% \end{minipage}
% 
% \section{Package macros}
% 
% This package has the following commands.
% 
% \begin{function}{\subtitle}
%   \begin{syntax}
%     \cs{subtitle}\marg{a subtitle}
%   \end{syntax}
%   The subtitle is printed exactly under the title on the
%   title-page.
% \end{function}
%
% \begin{function}{\shorttitle}
%   \begin{syntax}
%     \cs{shorttitle}\marg{a short title}
%   \end{syntax}
%   The short title is printed on the top right corner of
%   the page with the year of submission.
% \end{function}
% 
% \begin{function}{\supervisor}
%   \begin{syntax}
%     \cs{supervisor}\marg{name of the supervisor}
%   \end{syntax}
%   Name of the supervisor is printed on the title-page, in
%   the declaration text and under the name of the
%   supervisor for the signature.
% \end{function}
%
% \section{Package options}
% 
% \begin{function}{digital}
%   This package option is developed for getting a read-only
%   version of the thesis. By default the thesis is
%   formatted as a print-ready document as the per the
%   university requirements. If you want to submit a digital
%   copy of the thesis, use this option.
% \end{function}
% 
% \begin{function}{lof}
%   This package option is developed for getting a list of
%   figures after the list of tables. Every thesis might not
%   need a list of figures, therefore it isn't added by
%   default.
% \end{function}
% 
% \section{Example code}
% 
% Try this sample document.
%
% \begin{verbatim}
% \documentclass[digital,lof]{muling}
% \title{Linguistic profile of Konkani}
% \author{Niranjan}
% \supervisor{Dr.\ Renuka Ozarkar}
% \shorttitle{Konkani}
% \subtitle{%
%   Phonological \& morphological description of the Konkani
%   language%
% }
% 
% \begin{document}
% \maketitle
% 
% \chapter{Test chapter}
% \pex
% Department of linguistics
% \xe
% 
% \begin{center}
%   \begin{forest}
%     [NP%
%       [\head{N}\\
%         Department%
%       ]%
%       [PP%
%         [\head{P}\\%
%           of%
%         ]%
%         [NP%
%           [\head{N}\\%
%             linguistics%
%           ]%
%         ]%
%       ]%
%     ]%
%   \end{forest}
% \end{center}
% \end{document}
% \end{verbatim}
% \end{documentation}
% 
% \begin{implementation}
% \section{Implementation}
%    \begin{macrocode}
%<*class>
%    \end{macrocode}
%    \begin{macrocode}
\ProvidesClass{muling}[2023/04/02 v0.5
Department of Linguistics, University of Mumbai]
\DeclareOption{digital}{\PassOptionsToClass{oneside}{book}}
\newif\iflof
\DeclareOption{lof}{\loftrue}
\ProcessOptions
\iflof
\def\@printlof{%
  \thispagestyle{empty}%
  \phantomsection
  \addcontentsline{toc}{section}{List of figures}%
  \listoffigures
  \clearpage\pagebreak
}%
\else
\let\@printlof\@empty
\fi
\LoadClass{book}
%    \end{macrocode}
% This part of code loads the \cls{book} class \textit{\&}
% provides the option \verb|digital| for typesetting a
% document with read-only format \textit{\&} option
% \verb|lof| for automatically adding list of
% figures. \href{https://topanswers.xyz/tex?q=1500#a1728}{Phelype's
% answer} helped me for developing this code.
%    \begin{macrocode}
\let\@supervisor\@empty
\def\supervisor#1{%
  \def\@supervisor{#1}%
}
\let\@shorttitle\@empty
\def\shorttitle#1{%
  \def\@shorttitle{#1}%
}
\let\@subtitle\@empty
\def\subtitle#1{%
  \def\@subtitle{#1}%
}
%    \end{macrocode}
% This code enables three additional entries required in the
% title-page \& in the declaration \textit{\&} at various
% other places.
%    \begin{macrocode}
\RequirePackage{csquotes}
\RequirePackage[backend=biber,style=apa]{biblatex}
\RequirePackage[linguistics]{forest}
\RequirePackage{tipa}
\RequirePackage{leipzig}
\RequirePackage{expex}
\RequirePackage{fancyhdr}
\RequirePackage{fontspec}
\RequirePackage{hyperref}%
\makeglossaries
\setmainfont[%
  % Courtesy: https://topanswers.xyz/tex?q=3042#a3083
  IgnoreFontspecFile,%
  StylisticSet      = {5},%
  Renderer          = {Harfbuzz},%
  ItalicFont        = {NewCM10-Italic.otf},%
  BoldFont          = {NewCM10-Bold.otf},%
  BoldItalicFont    = {NewCM10-BoldItalic.otf},%
  SmallCapsFeatures = {%
    % Courtesy: https://tex.stackexchange.com/a/107246
    Numbers         =   {Lining,Monospaced},%
    Letters         =   {SmallCaps}%
  }%
]{NewCM10-Regular.otf}
\pagestyle{fancy}
\cfoot{\thepage}
\rhead{}
\lhead{%
  \begin{tabular}{l}%
    \@author\\%
    \textsc{%
    \footnotesize University of Mumbai%
    }%
  \end{tabular}%
  \hfill
  \@shorttitle, \the\year
}%
\AddToHook{begindocument/before}{%
  \colorlet{mulingred}{red!60!black}
  \colorlet{mulinggreen}{green!60!black}
  \colorlet{mulingblue}{blue!60!black}
  \hypersetup{%
    unicode,%
    colorlinks,%
    linkcolor     = {mulingred},%
    citecolor     = {mulinggreen},%
    urlcolor      = {mulingblue},%
    pdftitle      = {\@title},%
    pdfauthor     = {\@author},%
    pdfcreator    = {\@author},%
    pdfsubject    = {\@subtitle},%
    pdfkeywords   = {\@shorttitle}%
  }%
}
%    \end{macrocode}
% This part of code loads some other essential packages for
% typesetting documents with linguistic content \textit{\&}
% the name of the author, the name of the university and a
% short-title with the year of submission at the top right
% corner of every page.
% \href{https://tex.stackexchange.com/a/525452/174620}^^A
% {This} answer helped me for developing this code.
%    \begin{macrocode}
\newlength{\shift}
\setlength{\shift}{0.3in}
%    \end{macrocode}
% A new length is created for typesetting the title-page at
% the center of the page like in the \cls{article}
% class. \href{https://tex.stackexchange.com/a/514536/174620}%
% {This} answer helped me for developing this code.
%    \begin{macrocode}
\renewcommand{\maketitle}{%
  \thispagestyle{empty}%
  \addtolength{\oddsidemargin}{\shift}%
  \fboxsep3em
  \noindent\fbox{%
    \begin{minipage}%
      [c][\dimexpr\textheight-2\fboxsep-2\fboxrule]%
      [c]{\dimexpr\linewidth-2\fboxsep-2\fboxrule}%
      
      \begin{center}%
        \bigskip
        \hrule
        \vspace{2ex}%
        \textbf{%
          {\LARGE \@title}\\%
          \vspace{0.2cm}%
          {\textsc\@subtitle}%
        }%
        \vspace{2ex}%
        \hrule
        \vfill
        {%
          \large
          \textsc{%
            \@author
          }%
        }\\%
        \vspace{0.5in}%
        {%
          \large
          \textsc{%
            Supervisor: \@supervisor
          }%
        }\\%
        \vspace{0.5in}%
        {%
          \large
          \textsc{%
            A dissertation submitted in\\ the partial
            fulfillment of the requirements for the masters
            of linguistics%
          }%
        }\\%
        \vspace{0.5in}%
        \textsc{%
          Department of Linguistics\\%
          University of Mumbai, Kalina%
        }\\%
        \smallskip
        \bigskip
        \textsc{\@date}%
      \end{center}%
    \end{minipage}%
    \addtolength{\oddsidemargin}{-\shift}%
  }
  \clearpage\pagebreak
  \thispagestyle{empty}%
  \tableofcontents
  \clearpage\pagebreak
  \thispagestyle{empty}%
  \phantomsection
  \addcontentsline{toc}{section}{Abbreviations}%
  \printglossary[style=mcolblock]
  \clearpage\pagebreak
  \thispagestyle{empty}%
  \phantomsection
  \addcontentsline{toc}{section}{List of tables}%
  \listoftables
  \clearpage\pagebreak
  \thispagestyle{empty}%
  \phantomsection
  \begin{center}%
    {%
      \LARGE
      Declaration%
    }
  \end{center}%

  \bigskip
  As required by the University Regulation no.: R.\ 1972 I
  wish to state that the work embodied in this thesis
  titled \enquote{\@title: \@subtitle{}} forms my own
  contribution to the research work carried out under the
  guidance of \@supervisor{} at the University of Mumbai.
  
  This work has not been submitted for any other
  degree of this or any other university. Whenever
  references have been made to previous works of others,
  it has been clearly indicated as such and included in
  the bibliography.

  \bigskip
  \flushright{%
    Date: \@date\\%
    \vspace{0.5in}%
  }%
  
  \noindent Signature of the candidate:\\%
  \hfill Name: \@author
  \\%

  \bigskip
  \flushleft{%
    Certified by:\\%
    \vspace{0.5in}%
    Signature of the guide:\\%
    Name: \@supervisor
  }%
}%
%    \end{macrocode}
% This code typesets a customized title-page, adds a table
% of contents, list of tables\footnote{Having a list of
% figures is optional as seen earlier.} and the declaration
% text with all the details filled \textit{\&} some blank
% space left for the signature.
%    \begin{macrocode}
%</class>
%    \end{macrocode}
% \end{implementation}
% \Finale